\section{Constructing Machine-Checked Proofs}
\label{sec:proofs}

The core proof pattern in this project is semantic rather than syntactic.
Instead of proving correctness of a global optimizer function, we state a
refinement theorem directly between two definitions and unfold the semantic
interpretation enough to show that whenever the target program on the
left-hand side returns some value, the source program on the right-hand side
can either return that same value exactly or justify it through the
poison-aware register-value refinement relation. In most proofs in this
repository, the exact-equality branch is enough.

\subsection{A Simple Refinement: \texttt{add x, 0}}

The simplest examples are algebraic identities such as:
\small
\begin{leanlisting}
theorem ret_add_x_0_is_refined_by_ret_x :
  [llvm-definition|
    define i8 @f(i8 %x) {
    entry:
      ret i8 %x
    }
  ] ⊑ [llvm-definition|
    define i8 @f(i8 %x) {
    entry:
      %x = add i8 %x, 0
      ret i8 %x
    }
  ] := by
  ...
\end{leanlisting}
\normalsize

The shape of the proof is straightforward but still instructive. We first
introduce the refinement assumptions: arguments, return value,
well-formedness of both definitions, signature compatibility, and
well-formedness of argument values. Then we case-split on the input argument
list to discard impossible shapes. Finally, the actual semantic goal reduces
to a bitvector fact that simplification can solve.

This style scales well to a width-generic variant because width instantiation
is explicit in the syntax and semantics. The theorem
\texttt{ret\_add\_x\_0\_is\_refined\_by\_ret\_x\_generic} shows that the same proof
idea survives after replacing a fixed width by a symbolic width parameter.
Here \texttt{llvm-1-definition} is a quasiquoter for definitions with one
symbolic width parameter, later instantiated by \texttt{singletonWidths w}.
\small
\begin{leanlisting}
theorem ret_add_x_0_is_refined_by_ret_x_generic (w : Nat) :
  [llvm-1-definition|
    define i$0 @f(i$0 %x) {
    entry:
      ret i$0 %x
    }
  ].instantiateWidths (singletonWidths w)
  ⊑
  [llvm-1-definition|
    define i$0 @f(i$0 %x) {
    entry:
      %x = add i$0 %x, 0
      ret i$0 %x
    }
  ].instantiateWidths (singletonWidths w) := by
  ...
\end{leanlisting}
\normalsize

\subsection{A More Interesting Refinement: \texttt{undef}}

The more distinctive example is the theorem relating
\texttt{add undef, undef} and \texttt{mul undef, 2}. In the current project
this theorem is stated for every instantiated width:
\small
\begin{leanlisting}
theorem undef_add_is_refined_by_undef_mul2_generic (w : Nat) :
  [llvm-1-definition|
    define i$0 @f() {
      entry:
        %x = mul i$0 undef, 2
        ret i$0 %x
    }
  ].instantiateWidths (singletonWidths w)
  ⊑
  [llvm-1-definition|
    define i$0 @f() {
      entry:
        %x = add i$0 undef, undef
        ret i$0 %x
    }
  ].instantiateWidths (singletonWidths w) := by
  ...
\end{leanlisting}
\normalsize

What makes this theorem interesting is that it is not simply an algebraic
identity on the same input value. The two occurrences of \vundef on the right
may resolve independently, while the one occurrence on the left resolves to a
single concrete value. Since the left-hand side is the target and the
right-hand side is the source, the proof uses the refinement definition's
existential freedom on the right-hand-side (source) supply. Given a supply for the
multiplication program, the proof constructs a matching supply for the
addition program by repeating the observed choice. The arithmetic lemma then
shows that doubling a bitvector by addition yields the same result as
multiplication by two, so the proof can use the exact-equality branch of the
refinement disjunction.
\small
\begin{leanlisting}[numbers=left,xleftmargin=0.75em]
  intro args hwfAdd hwfMul hsig hargs ~\label{line:undef-proof-intro}~
  constructor
  · intro hdefined u' ~\label{line:undef-proof-definedness}~
    cases args with ~\label{line:undef-proof-args}~
    | nil => ~\label{line:undef-proof-nil-args}~
      have hnil : False := by
        have hrun : ... = Except.error "" := by ~\label{line:undef-proof-empty-supply}~
          simp
          rfl
        obtain ⟨ret_x, hret_x⟩ := hdefined []
        simp_all
      exact False.elim hnil
    | cons arg rest => ~\label{line:undef-proof-impossible-args}~
      simp [Definition.ArgValuesWellFormed] at hargs
  · intro u hRet h ~\label{line:undef-proof-value}~
    cases u with ~\label{line:undef-proof-supply-split}~
    | nil =>
        have hnil : ... = Except.error "" := by
          simp
          rfl
        simp_all
    | cons x xs => ~\label{line:undef-proof-cons-supply}~
      have hmulRun : ... =
          Except.ok (RegisterValue.bv w
            (.value ((BitVec.ofNat w x) * (2#w)))) := by ~\label{line:undef-proof-hmul}~
        simp [simp_llvm]
        rfl
      refine ⟨[x, x], Or.inl ?_⟩ ~\label{line:undef-proof-witness}~
      have haddRun : ... =
          Except.ok (RegisterValue.bv w
            (.value ((BitVec.ofNat w x) + (BitVec.ofNat w x)))) := by ~\label{line:undef-proof-hadd}~
        simp [simp_llvm]
        rfl
      simp_all [bitvec_double_eq_mul_two_generic] ~\label{line:undef-proof-simpall}~
\end{leanlisting}
\normalsize

The proof proceeds in four steps.
\begin{enumerate}
  \item At line~\ref{line:undef-proof-intro}, the script introduces the refinement
    hypotheses. The constructor immediately splits the goal into the
    definedness obligation at line~\ref{line:undef-proof-definedness} and the
    value-preservation obligation at line~\ref{line:undef-proof-value}. Within the
    definedness branch, the argument list is inspected at
    line~\ref{line:undef-proof-args}; since the function takes no arguments, the
    impossible nonempty case at line~\ref{line:undef-proof-impossible-args} is
    discharged from \texttt{Definition.ArgValuesWellFormed}.
  \item The definedness argument is resolved by contradiction in the empty
    argument branch at line~\ref{line:undef-proof-nil-args}. The auxiliary run fact
    \texttt{hrun} at line~\ref{line:undef-proof-empty-supply} shows that the source
    program \texttt{add undef, undef} fails on the empty supply. That
    contradicts the hypothesis \texttt{hdefined}, which claims success for
    every supply, so \texttt{simp\_all} closes the branch.
  \item The value-preservation branch begins at
    lines~\ref{line:undef-proof-value} and \ref{line:undef-proof-supply-split}. After the
    trivial empty-supply case, the interesting branch is the nonempty target
    supply \texttt{x :: xs} at line~\ref{line:undef-proof-cons-supply}. The fact
    \texttt{hmulRun} at line~\ref{line:undef-proof-hmul} computes the target run as
    multiplication by two. The witness chosen at
    line~\ref{line:undef-proof-witness} is the source supply \texttt{[x, x]},
    forcing both occurrences of \vundef to resolve to the same concrete value;
    the proof then enters the left branch of the refinement disjunction,
    corresponding to exact return equality.
    The fact \texttt{haddRun} at line~\ref{line:undef-proof-hadd} then computes the
    matching source result as \texttt{x + x}.
  \item At line~\ref{line:undef-proof-simpall}, \texttt{simp\_all} combines the
    execution hypothesis with the two concrete run facts and reduces the
    remaining goal to the arithmetic identity
    \texttt{BitVec.ofNat w x + BitVec.ofNat w x = BitVec.ofNat w x * 2}, which
    is discharged by the helper lemma below.
\end{enumerate}

For reference, the helper lemma used in the final step is:
\small
\begin{leanlisting}
private theorem bitvec_double_eq_mul_two_generic {w : Nat} (b : BitVec w) :
    b + b = b * (2 : BitVec w) := by
  ...
\end{leanlisting}
\normalsize
Its proof is routine bitvector arithmetic and is omitted from the report.

This example captures an important lesson from the project: once \vundef is
present, even very small optimization theorems stop being purely local
rewrites. The proof has to say something precise about how nondeterministic
choices line up between the two sides. Making that alignment explicit was one
of the main conceptual advances of the final stage of the project.

The generic story is also sharp in the opposite direction. For positive
widths, the converse refinement fails:
\small
\begin{leanlisting}
theorem undef_mul2_is_not_refined_by_undef_add_generic (w : Nat) (hpos : 0 < w) :
  ¬ ([llvm-1-definition|
    define i$0 @f() {
      entry:
        %x = add i$0 undef, undef
        ret i$0 %x
    }
  ].instantiateWidths (singletonWidths w)
  ⊑
  [llvm-1-definition|
    define i$0 @f() {
      entry:
        %x = mul i$0 undef, 2
        ret i$0 %x
    }
  ].instantiateWidths (singletonWidths w)) := by
  ...
\end{leanlisting}
\normalsize

The counterexample is simple. The addition program can return
\texttt{BitVec.ofNat w 1} using the supply \texttt{[1, 0]}, because the two
\vundef occurrences are chosen independently. Any execution of
\texttt{mul undef, 2}, however, returns a value of the form \texttt{b * 2}, and
the helper lemma \texttt{bitvec\_mul\_two\_ne\_one\_generic} shows that such a
result can never equal \texttt{1} when \texttt{0 < w}. So the asymmetry in the
refinement relation is not an artifact of a fixed-width example; it persists
uniformly across all positive bitwidths.
